\bookchapter{Appendice F: Speciale acquisto e uso degli schermi pieghevoli}

\chapterintro{Lo schermo pieghevole è passato da “tecnologia del futuro” a categoria matura “che si può comprare ora”. Ma rispetto a un telefono a barra differisce molto per forma, abitudini d'uso e manutenzione: vale la pena comprarlo, come sceglierlo, come usarlo — questo capitolo ti aiuta a fare chiarezza.}

\booksection{A che punto sono oggi gli schermi pieghevoli?}

Prima la conclusione:\textbf{lo schermo pieghevole non è più un “giocattolo da assaggiare”, ma non è nemmeno una forma mainstream “che tutti dovrebbero comprare”}. Sta a metà tra prodotto di nicchia per geek e bene di consumo di massa, adatto a specifici pubblici.

Attualmente si divide in due categorie principali:

\begin{itemize}
  \item \textbf{Pieghevole orizzontale (stile libro)}: chiuso è grande come un telefono normale, aperto è un piccolo tablet quasi quadrato. Rappresentativi: serie Huawei Mate X, serie Samsung Galaxy Z Fold, serie Xiaomi MIX Fold, serie vivo X Fold. Promuove "un telefono = telefono + tablet".
  \item \textbf{Pieghevole verticale (stile a conchiglia)}: chiuso sembra un portacipria, aperto è un telefono di dimensioni normali. Rappresentativi: serie Huawei Pocket, serie Samsung Galaxy Z Flip, OPPO Find N Flip, Xiaomi MIX Flip. Promuove "compatto e portabile + sguardo rapido su schermo esterno".
\end{itemize}

Progressi chiave del 2024-2026: durata della cerniera dichiarata comunemente 500.000 piegamenti (100 aperture/chiusure al giorno per 13 anni), vetro ultra-sottile UTG che sostituisce le pellicole plastiche iniziali, piega visibile a occhio nudo ma già poco percettibile al tatto, prezzo sceso dalla fascia 10.000 yuan a 4.000-5.000 yuan.

\begin{tcbinfo}{Giudizio in una frase}
Se oggi un telefono a barra ti basta completamente e non hai il punto dolente di “schermo non abbastanza grande”, non comprare il pieghevole. È un dispositivo che “aggiunge fiori al brodo”, non uno che “porta carboni al freddo”.

\end{tcbinfo}

\booksection{Chi è adatto al pieghevole?}

In base ai feedback reali degli utenti, queste categorie hanno un'alta probabilità di “innamorarsene”:

\begin{enumerate}
  \item \textbf{Lettori pesanti / elaboratori di documenti}: chi legge spesso PDF sul telefono, annota documenti, legge articoli lunghi. Lo schermo interno da quasi 7,5 pollici una volta aperto permette di leggere PDF A4 senza zoom, un salto di qualità.
  \item \textbf{Utenti business che viaggiano / vanno spesso a riunioni}: non vogliono portarsi il tablet ma trovano piccolo lo schermo del telefono. Il pieghevole orizzontale aperto per guardare PPT, firmare contratti, confrontare email su due schermi è molto più efficiente di un telefono a barra.
  \item \textbf{Creativi amanti della modalità "sospesa"}: piegato a 90° sul tavolo, schermo superiore per l'inquadratura e inferiore per lo scatto (selfie / time-lapse), schermo superiore per il video e inferiore per appunti, schermo superiore per la live e inferiore per i commenti — una forma che un telefono a barra non può offrire.
  \item \textbf{Utenti donne che cercano portabilità e grande schermo}: il pieghevole verticale chiuso entra in una piccola borsa, aperto è un telefono normale, lo schermo esterno funge anche da specchietto e mostra notifiche.
\end{enumerate}

\begin{tcbwarn}{Queste categorie non dovrebbero comprarlo}
\begin{itemize}
  \item Chi usa il telefono solo per brevi video, WeChat e telefonate — il pieghevole non apporta nulla, anzi è più pesante e spesso.
  \item Chi ha mani piccole o tendiniti ai polsi — i pieghevoli orizzontali aperti pesano comunemente 230 g+, reggerli con una mano stanca.
  \item Budget stretto, un telefono da usare per 4+ anni — il valore di rivendita del pieghevole è inferiore al telefono a barra e i costi di riparazione sono alti.
\end{itemize}

\end{tcbwarn}

\booksection{Punti chiave per l'acquisto: guarda questi 6 indicatori}

Non farti abbagliare dalla parola “pieghevole”, in fase di acquisto guarda soprattutto:

\begin{itemize}
  \item \textbf{Tipo di cerniera}: cerniera a goccia (piega meno profonda, vita più lunga, standard attuale) migliore di quella a U (piega più profonda, soluzione iniziale). Chiedi al venditore o verifica nelle recensioni quale sia.
  \item \textbf{Materiale schermo interno}: deve essere UTG ultra-thin glass, non le pellicole plastiche iniziali. UTG al tatto è simile al vetro, la pellicola plastica sotto l'unghia avvallamenti e si segna facilmente.
  \item \textbf{Comportamento della piega}: vai in negozio e prova ad aprire di persona. Il livello accettabile nel 2026 è: a schermo acceso frontalmente quasi invisibile, in luce laterale si vede una leggera scanalatura, scorrendo col dito c'è un leggero avvallamento ma non taglia.
  \item \textbf{Dimensioni schermo esterno}: per i pieghevoli orizzontali più grande è meglio (meglio sopra i 6,4 pollici), chiuso si può usare come telefono normale; per i verticali più grande è più utile (sopra i 4 pollici per rispondere ai messaggi, vedere la navigazione).
  \item \textbf{Peso e spessore}: per i pieghevoli orizzontali aperti 230-260 g è la norma, sotto i 230 g è eccellente; spessore da chiuso 12-14 mm, sotto i 12 mm è sottile. I verticali chiusi devono entrare facilmente nella tasca della camicia per essere considerati accettabili.
  \item \textbf{Post-vendita e garanzia}: guarda soprattutto il prezzo di riparazione per lo schermo interno rotto. Alcuni marchi offrono la prima sostituzione gratuita nel primo anno, altri richiedono 2.000-4.000 yuan a proprie spese. È il costo nascosto maggiore del pieghevole.
\end{itemize}

\begin{tcbpractice}{Checklist per la prova in negozio}
\begin{enumerate}
  \item Apri e chiudi 10 volte con le tue mani, senti se lo smorzamento della cerniera è fluido e se ci sono rumori anomali.
  \item Aperto, guarda la piega di taglio verso la luce, poi frontalmente a schermo acceso.
  \item Tieni aperto con una mano per 30 secondi, valuta se il peso è accettabile.
  \item Prova la modalità sospesa: piega a 90° sul tavolo, prova a scattare foto e guardare video.
  \item Chiedi chiaramente il prezzo di riparazione sola schermo interno e la politica di garanzia.
\end{enumerate}

\end{tcbpractice}

\booksection{Trucchi d'uso unici del pieghevole}

Dopo averlo comprato, queste sono cose che un telefono a barra non può offrire:

\booksubsection{Split-screen e multi-tasking}

Aperto, la maggior parte dei sistemi supporta split-screen sopra/sotto o sinistra/destra. Scenari tipici: schermo superiore per video-lezione, inferiore per appunti; sinistra WeChat destra documento; sopra mappe sotto musica. Molto più utile dello split-screen di un telefono a barra, perché lo schermo è abbastanza grande e ciascuna metà non è stipata.

\booksubsection{Modalità sospesa (piegato a 90°)}

È l'uso “killer” del pieghevole:

\begin{itemize}
  \item \textbf{Selfie / foto di gruppo}: posato sul tavolo piegato a 90°, schermo superiore per l'inquadratura, inferiore per lo scatto, senza dover chiedere aiuto.
  \item \textbf{Time-lapse / notturna}: stesso, posato sul tavolo per riprese stabili, risparmia il treppiede.
  \item \textbf{Videochiamata}: mani libere, l'angolo con cui l'altro ti vede è più naturale.
  \item \textbf{Guardare video}: posato sul tavolo come una piccola TV, particolarmente utile sui vassoi di aereoo/treno ad alta velocità.
  \item \textbf{Anteprima schermo esterno + ripresa schermo interno}: quando usi la fotocamera posteriore per riprendere te stesso, lo schermo esterno mostra in tempo reale l'immagine, la composizione è più precisa.
\end{itemize}

\booksubsection{Sguardo rapido schermo esterno}

Il pieghevole verticale chiuso, con lo schermo esterno, può: vedere ora/notifiche/meteo, rispondere a WeChat (voce-testo o risposte rapide), pagare i mezzi pubblici, vedere la navigazione, fare da specchietto (fotocamera anteriore). Senza doverlo aprire ogni volta, risparmia batteria e fatica.

\booksubsection{Modalità computer}

Alcuni pieghevoli orizzontali (Huawei, Samsung) aperti possono passare in “modalità computer”; con monitor e tastiera/mouse esterni diventano un PC elementare, sufficiente per elaborare documenti in emergenza e rispondere alle email. Non è uno strumento di produttività principale, ma in emergenza è molto comodo.

\booksection{Manutenzione: far vivere più a lungo il pieghevole}

La parte più delicata del pieghevole è lo schermo interno e la cerniera; queste abitudini possono prolungarne sensibilmente la vita:

\begin{itemize}
  \item \textbf{Non forzare l'apertura a basse temperature}: sotto lo 0 °C il vetro UTG diventa fragile, prima di aprirlo scaldalo qualche minuto in tasca. È un punto che tutti i manuali scrivono ma nessuno legge.
  \item \textbf{Evita il contatto dello schermo interno con unghie o oggetti duri}: lo schermo interno ha uno strato protettivo ma è più morbido dell'esterno; usa il polpastrello, non le unghie.
  \item \textbf{Prima di chiudere verifica che non ci siano corpi estranei}: un granello di sabbia in mezzo alla chiusura può rompere lo schermo interno. Dai un'occhiata prima di chiudere.
  \item \textbf{Non applicare pellicole temperate comuni}: lo schermo interno va protetto solo con pellicole ufficiali dedicate o pellicole UTG; le pellicole temperate comuni possono compromettere piega e touch. Si raccomanda di far applicare la pellicola presso l'assistenza ufficiale.
  \item \textbf{Non far entrare polvere o liquidi nella cerniera}: anche se molti sono dichiarati IPX8, il grado anti-polvere è comunemente basso. Attenzione in spiaggia, cantieri e simili.
  \item \textbf{Forza di apertura uniforme}: non aprire con uno strappo violento a una mano; apertura a due mani dal centro, uniforme, è la più stabile.
\end{itemize}

\textbf{Scenari ad alto rischio di danneggiamento dello schermo interno}\par

\begin{itemize}
  \item Chiudere tenendo dentro il cavo delle cuffia / di ricarica — fessura diretta dello schermo interno.
  \item In borsa mescolati con le chiavi — lo schermo esterno non ha problemi, quello interno viene incrinato dalle chiavi.
  \item In inverno, aprire e chiudere subito dopo essere entrati dal freddo in una stanza calda — il salto termico altera le tensioni del vetro.
  \item Dopo un anno, pensare che "la piega si è approfondita" e staccare la pellicola originale — staccandola si può danneggiare lo schermo interno.
\end{itemize}

\booksection{Le “trappole” del pieghevole e aspettative mentali}

Prima di comprare, devi fare i conti con questi punti:

\begin{itemize}
  \item \textbf{La piega ci sarà sempre}: la tecnologia del 2026 non può "azzerare la piega", può solo far sì che "a schermo acceso frontalmente quasi non si veda". Se sei un perfezionista che ogni giorno soffre a guardare la piega, non comprare.
  \item \textbf{Più pesante e spesso}: i pieghevoli orizzontali aperti pesano comunemente 230 g+, reggerli con una mano stanca più di un telefono a barra; chiusi sono spessi 12-14 mm, più gonfi in tasca.
  \item \textbf{Autonomia comunemente media}: dovendo ospitare il meccanismo della cerniera, lo spazio per la batteria è compresso; a parità di prezzo l'autonomia è di solito inferiore a un telefono a barra.
  \item \textbf{Valore di rivendita basso}: dopo un anno il prezzo di usato è di solito il 50-60\% del nuovo, scende più velocemente del telefono a barra.
  \item \textbf{Riparazioni care}: riparazione schermo interno 2.000-4.000 yuan, riparazione cerniera 500-1.500 yuan. Prima di comprare verifica la politica di garanzia.
  \item \textbf{Adattamento app disomogeneo}: le app principali si sono adattate al grande schermo, ma le app di nicchia possono solo "essere allungate", con esperienza media.
\end{itemize}

Il pieghevole non è un “telefono migliore”, ma “un'altra forma di telefono”. Lo compri perché hai bisogno della sua forma unica, non perché è “più avanzato”.

— Filosofia di acquisto dei pieghevoli

\booksection{Come scegliere tra i principali marchi (2026)}

Non scrivo modelli specifici (diventerebbero superati), solo il carattere del marchio per aiutarti a orientarti rapidamente:

\begin{itemize}
  \item \textbf{Serie Huawei Mate X}: punto di riferimento del pieghevole orizzontale, forte integrazione dell'ecosistema HarmonyOS, le proporzioni dello schermo esterno sono vicine a un telefono normale, chiuso si può usare come telefono principale. Prezzo alto ma valore di rivendita relativamente migliore.
  \item \textbf{Samsung Galaxy Z Fold/Flip}: pioniere del pieghevole, taratura della cerniera matura, Z Flip è il punto di riferimento estetico del verticale. Sul territorio nazionale la rete di assistenza è meno diffusa.
  \item \textbf{Xiaomi MIX Fold/Flip}: linea di rapporto qualità-prezzo, dotazioni ricche, prezzo relativamente accessibile. Adatto a chi vuole provare senza spendere troppo.
  \item \textbf{Serie vivo X Fold}: tra i pieghevoli orizzontali, le proporzioni dello schermo esterno sono le più vicine a un telefono normale, chiuso si usa bene come telefono a barra.
  \item \textbf{Serie OPPO Find N}: tra i pieghevoli orizzontali, dimensioni relativamente compatte, si tiene con una mano, adatto a chi trova i grandi pieghevoli troppo pesanti.
  \item \textbf{Serie Huawei Pocket}: tra i verticali, lo schermo esterno è il più grande e utile, le funzioni orientate al femminile come specchietto e rilevamento UV sono curate.
\end{itemize}

Consigli d'acquisto: budget ampio e priorità all'ecosistema scegli Huawei Mate X; per la cerniera più matura scegli Samsung Z Fold; per rapporto qualità-prezzo Xiaomi/vivo; per il verticale Huawei Pocket e Samsung Z Flip.

\begin{tcbinfo}{Ultimo avvertimento}
La tecnologia pieghevole itera ogni anno; questa appendice è scritta nel 2026. Prima di comprare consulta le ultime recensioni e i feedback degli utenti, con attenzione a tre punti: test reali della durata della cerniera, comportamento della piega, politica di riparazione dello schermo interno. Toccarlo di persona in negozio conta sempre più delle recensioni.

\end{tcbinfo}
